\begin{textsample}{POS Dim 2 – 2012\_c6 – Score -1.00 – 2012\_c6/enade\_2012\_dir\_5ad34d9588.txt}  \label{ex:f2_pos_248}
CAPÍTULO 3 ANÁLISE TÉCNICA DA PROVA Este capítulo tem por objetivo apresentar o desempenho dos estudantes concluintes de Direito no ENADE/2012.
Para isso, foram calculadas as estatísticas básicas da prova em seu todo, bem como as estatísticas dos componentes relacionadas à Formação Geral, ao de Conhecimento Específico da Área e das questões discursivas isoladamente.
Nas tabelas, são apresentados o tamanho da população inscrita e de presentes, e as seguintes estatísticas das notas9 : média do desempenho na prova, erro padrão da média, desvio padrão, nota mínima, mediana e nota máxima.
Tais estatísticas contemplam o total de estudantes concluintes da Área de Direito presentes à prova do ENADE 2012, tendo -se em vista as seguintes agregações : ( a ) as Grandes Regiões e o país como um todo ; ( b ) a Categoria Administrativa ; e ( c ) a Organização Acadêmica.
Em relação aos gráficos de distribuição de notas, o intervalo considerado foi de 10 unidades, aberto à esquerda e fechado à direita, com exceção do primeiro intervalo, [ 0 ; 10 ], fechado em ambos os extremos.
Para os gráficos de distribuição das notas das questões discursivas, foram consideradas mais duas categorias : questão em branco e nota zero. 3.1 ESTATÍSTICAS BÁSICAS DA PROVA 3.1.1 Estatísticas Básicas Gerais A Tabela 3.1 apresenta as estatísticas básicas da prova por grande Região.
A população total de inscritos foi de 123.211.
Destes, 104.662 estiveram presentes, sendo 15,1 \% o índice de não comparecimento.
A Região de maior abstenção foi a Centro-Oeste ( 17,2 \% ) e a de menor abstenção foi a Sul ( 12,2 \% ).
A média das notas da prova como um todo ( nas seções seguintes serão analisados os componentes de Formação Geral e de Conhecimento Específico ) foi 40,8, sendo que os alunos da região Centro-Oeste obtiveram a média mais baixa ( 40,0 ) e os da região Sul obtiveram a média mais alta ( 41,5 ).
As demais médias foram : 40,4 na região Norte, 40,3 na região Nordeste e 41,1 na região Sudeste.
O desvio padrão para o Brasil como um todo foi 13,2, sendo o maior desvio padrão encontrado na região Nordeste ( 13,9 ) e o menor na região Norte ( 12,2 ), indicando uma menor dispersão das notas desta última região. 9 Essas estatísticas e outras estão definidas no Capítulo 1.
A região que obteve a maior nota máxima foi a Centro-Oeste ( 88,0 ), ao passo que a região que atingiu a menor nota máxima foi a Norte ( 79,4 ).
A mediana do Brasil como um todo foi 41,2, sendo a maior mediana obtida na região Sul ( 41,9 ) e a menor obtida nas regiões Norte e Centro-Oeste ( 40,3 ).
A nota mínima foi zero em todas as regiões, sem exceção.
O comportamento das notas dos estudantes de todo o Brasil pode ser observado no Gráfico 3.1 que apresenta um histograma com a distribuição das mesmas.
Essa é uma distribuição unimodal com moda no intervalo ( 40 ; 50 ].
Apesar do coeficiente de assimetria da distribuição das notas ser negativo ( -0,17 ) este é pequeno e podemos considerar que a distribuição é aproximadamente simétrica.
As distribuições por Grande Região também apresentam assimetria negativa ( variando de -0,23 até -0,02 ), concentração pouco maior do lado direito do histograma e mais espalhada do lado esquerdo.
Os Gráficos 3.2, 3.3 e 3.4 apresentam informações referentes à média da nota final dos Participantes, desagregando os resultados de acordo com, respectivamente, as Grandes Regiões do país, a Categoria Administrativa e a Organização Acadêmica.
Os gráficos apresentam o valor da média das notas como uma barra e os extremos do intervalo de confiança de 95 \% como linhas verticais unidas por uma linha horizontal na forma da letra H maiúscula.
Considerando -se o gráfico de notas segundo Grande Região ( Gráfico 3.2 ), observa -se que existe diferença estatisticamente significativa ao nível de 95 \% entre a maior média, obtida na região Sul ( 41,5 ) e a menor, obtida na região Centro-Oeste ( 40,0 ).
Levando -se em conta os estudantes por Categorias Administrativas ( Gráfico 3.3 ), observa -se que existe diferença estatisticamente significante entre as médias das notas das IES Públicas e Privadas.
Pode ser observado ainda que a média dos alunos de IES Públicas é maior ( 44,6 ) que a dos alunos de IES Privadas ( 40,3 ).
A diferença entre IES Públicas e Privadas ( 4,3 ) é maior do que as diferenças entre regiões Sul e Centro-Oeste ( 1,5 ), a maior e a menor média.
Tendo como foco o Gráfico 3.4, que apresenta as notas médias das provas segundo Organização Acadêmica, constata -se que existe diferença estatisticamente significativa ao nível de 95 \% nas médias das notas dos estudantes provenientes de Universidades, Centros Universitários e Faculdades.
A maior média foi obtida pelos estudantes de Universidades ( 41,7 ), e a menor, pelos de Faculdades ( 40,0 ). 3.1.2 Estatísticas Básicas no Componente de Formação Geral A Tabela 3.2 apresenta as estatísticas básicas em relação ao componente da prova que avalia a Formação Geral dos estudantes concluintes.
Os alunos de todo Brasil obtiveram desempenho médio de 46,0.
Quanto à variabilidade, o desvio padrão das notas dos estudantes do Brasil como um todo foi 17,4.
A maior média foi obtida na região Sudeste ( 46,4 ), e a menor, na região Nordeste ( 45,4 ).
As demais médias foram : 46,2 na região Norte, 46,0 na região Sul e 45,6 na região Centro-Oeste.
Já o maior desvio padrão foi obtido na região Nordeste ( 18,4 ) e o menor na região Norte ( 16,6 ).
Os demais desvios padrões foram : 17,5 na região Centro-Oeste, 17,1 na região Sudeste e 16,9 na região Sul.
A maior nota no componente de Formação Geral da prova do ENADE foi obtida por pelo menos um aluno da região Sudeste ( 96,0 ) enquanto que a menor nota máxima foi obtida na região Norte ( 88,5 ).
Nas outras regiões as notas máximas foram as mesmas, 95,0 nas regiões Nordeste, Sul e Centro-Oeste.
A mediana do Brasil como um todo foi 47,0, sendo a menor mediana encontrada nas regiões Nordeste e Centro-Oeste ( 46,5 ) e a maior encontrada na região Sudeste ( 47,5 ).
A nota mínima nesta parte foi zero em todas as regiões, sem exceção.
O Gráfico 3.5 propicia a avaliação do desempenho dos estudantes no componente de Formação Geral a partir do histograma da distribuição das notas correspondentes.
A distribuição é unimodal, com moda em ( 40 ; 50 ], assim como na prova como um todo.
Nota -se, ainda, que no gráfico 3.5 as notas apresentam uma maior dispersão do que no Gráfico 3.1 ( distribuição das notas da prova ), confirmado pela comparação dos desvios padrões : 13,2 para a nota da prova como um todo e 17,4 para o componente de Formação Geral.
Destaca -se ainda, o intervalo ( 50 ; 60 ] que possui concentração muito próxima da modal.
Para o componente de Formação Geral, o coeficiente de assimetria da distribuição das notas dos estudantes, como na prova como um todo, também é negativo ( -0,33 ), porém maior em módulo, indicando um maior grau de assimetria.
Por ser um coeficiente pequeno, a distribuição apresenta uma leve concentração à direita e cauda maior à esquerda.
Em todas as Grandes Regiões os histogramas também possuem assimetria levemente negativa ( entre -0,35 para a região Nordeste, e -0,24 para a região Norte ).
Nos Gráficos 3.6, 3.7 e 3.8 são apresentadas as informações referentes ao desempenho dos Concluintes no componente de Formação Geral, em diferentes agregações : Grande Região do país, Categoria Administrativa e Organização Acadêmica.
Observa -se pelo Gráfico 3.6 que existe diferença estatisticamente significativa entre a maior média das notas no Componente de Formação Geral, obtida na região Sudeste ( 46,4 ), e a menor, obtida na região Nordeste ( 45,4 ).
Vemos que o intervalo de confiança mais largo é o da região Norte ; já o intervalo mais estreito é observado nas regiões Sudeste e Sul.
Este fato está relacionado, também, com o tamanho da população envolvida, menor na região Norte do que nas regiões Sudeste e Sul.
No Gráfico 3.7, que representa as notas médias no Componente de Formação Geral segundo Categoria Administrativa do país, observa -se que existe diferença estatisticamente significativa entre as médias.
Assim como ocorreu para as notas da prova como um todo, para o componente de Formação Geral, os concluintes de Direito, das IES Públicas ( 48,8 ) obtiveram uma média maior do que os das IES Privadas ( 45,7 ).
Considerando -se o tipo de Organização Acadêmica, nota -se, no Gráfico 3.8, uma diferença estatisticamente significativa ao nível de 95 \% nas médias das notas dos estudantes provenientes de Universidades, em relação aos provenientes de Centros Universitários e Faculdades.
Nas Universidades ( 46,7 ), a média é mais elevada do que em Centros Universitários ( 45,5 ) e Faculdades ( 45,4 ). 3.1.3 Estatísticas Básicas do Componente de Conhecimento Específico A Tabela 3.3 apresenta as estatísticas básicas referentes ao componente de Conhecimento Específico da área de Direito.
A média do desempenho dos alunos do Brasil como um todo foi 39,1.
A maior média foi obtida na região Sul ( 39,9 ), e a menor, na região Centro-Oeste ( 38,1 ).
As demais médias foram : 39,3 na região Sudeste, 38,5 na região Nordeste e 38,4 na região Norte.
Quanto à variabilidade das notas, o desvio padrão do Brasil como um todo foi 14,1, sendo o maior desvio padrão observado na região Nordeste ( 14,6 ) e o menor na região Norte ( 13,1 ).
Os demais desvios foram : 14,1 na região Centro-Oeste, 14,0 da região Sudeste e 13,9 da região Sul.
A mediana das notas dos estudantes de todo o Brasil foi 39,0.
A maior mediana ocorreu na região Sul ( 40,0 ) e a menor nas regiões Norte e Centro-Oeste ( 38,3 ).
As demais medianas foram 39,3 na região Sudeste e 38,5 na região Nordeste.
A nota máxima do Brasil como um todo foi 94,0, sendo obtida por pelo menos um aluno da região Centro-Oeste.
As demais notas máximas foram : 90,7 na região Nordeste, 88,8 na região Sudeste, 83,0 na região região Sul e 38,4 na região Norte.
A nota mínima foi zero em todas as regiões, sem exceção.
Assim como os Gráficos 3.1 e 3.5, o Gráfico 3.9, apresentado a seguir, proporciona uma avaliação do desempenho de concluintes em relação ao componente de Conhecimento Específico com um histograma da distribuição das notas correspondentes.
Dentre as 3 distribuições apresentadas, esta é a mais concentrada nas notas baixas.
Esta também é uma distribuição unimodal, e o grupo modal é o ( 30 ; 40 ], abaixo do grupo modal para a Formação Geral.
O coeficiente de assimetria da distribuição das notas do componente de Conhecimento Específico é negativo e bem próximo de zero ( -0,03 ).
Nota -se pelo histograma ( Gráfico 3.9 ) que esta é uma distribuição aproximadamente simétrica.
Já nas regiões Norte e Centro-Oeste, os coeficientes de assimetria são positivos, iguais a 0,10 e 0,05 respectivamente, evidenciando que a distribuição das notas dos estudantes dessas regiões, no componente de Conhecimento Específico, tem cauda ligeiramente mais pesada à direita.
Os Gráficos 3.10, 3.11 e 3.12 apresentam, respectivamente, uma comparação dos resultados em relação à Grande Região do país, à Categoria Administrativa e à Organização Acadêmica, agora levando em conta o desempenho de alunos no componente de Conhecimento Específico da prova.
Pelo Gráfico 3.10, observa -se que existe diferença estatisticamente significativa entre a maior média das notas no componente de Conhecimento Específico, obtida na região Sul ( 39,9 ), e a menor, obtida na região Centro-Oeste ( 38,1 ).
Quanto à Categoria Administrativa ( gráfico 3.11 ), observa -se um comportamento semelhante àquele da prova como um todo, ou seja, existe diferença estatisticamente significativa entre as médias das IES Públicas e Privadas, sendo que a maior média foi obtida por alunos de IES Públicas de ensino ( 43,1 ).
Quanto ao Gráfico 3.12, observa -se, mais uma vez, que existe diferença estatisticamente significativa ao nível de 95 \% entre as notas no componente de Conhecimento Específico dos diferentes tipos de Organização Acadêmica.
Sendo que a média dos Concluintes das Universidades ( 40,0 ) foi maior do que de Centros Universitários ( 38,5 ) e de Faculdades ( 38,2 ). 3.2 ANÁLISE DAS QUESTÕES OBJETIVAS 3.2.1 Componente de Formação Geral A Tabela 3.4 apresenta as estatísticas básicas relativas às oito questões objetivas do componente da prova que abrange a Formação Geral dos estudantes.
A média do Brasil foi 49,3.
A menor média foi encontrada na região Norte ( 48,7 ) e a maior na região Centro-Oeste ( 49,6 ).
As demais médias foram 49,3 na região Nordeste e 49,2 nas regiões Sudeste e Sul.
O desvio padrão do Brasil foi 21,1, sendo o maior desvio padrão encontrado na região Nordeste ( 21,8 ) e o menor na região Norte ( 20,4 ).
Os demais desvios foram : 21,3 na região Centro-Oeste, 21,0 na região Sudeste e 20,7 na região Sul.
As medianas ( 50,0 ), as notas máximas ( 100,0 ) e as notas mínimas ( 0,0 ), nas questões objetivas do componente de Formação Geral, foram iguais para todas as regiões, sem exceção.
A Tabela 3.5 apresenta os índices de facilidade e discriminação ( ponto bisserial ) das questões objetivas do componente de Formação Geral.
Quanto ao índice de facilidade, foram usadas as seguintes cores para diferenciar o nível de dificuldade da questão : • Azul para as questões classificadas com índice muito fácil ( índice \&gt ; =0,86 ), verde para as questões classificadas com índice fácil ( 0,61 a 0,85 ), amarelo para as questões classificadas com médio ( 0,41 a 0,60 ), vermelho para as questões classificadas com difícil ( 0,16 a 0,40 ) e roxo para as questões classificadas com muito difícil ( \&lt ; =0,15 ).
Já quanto ao índice de discriminação, foram usadas as seguintes cores para qualificar a questão : • As questões classificadas com índice fraco receberam a cor roxa ( índice \&lt ; =0,19 ), as classificadas com médio receberam a cor vermelho ( 0,20 a 0,29 ), as classificadas com bom receberam a cor verde ( 0,30 a 0,39 ) e as classificadas com muito bom ( \&gt ; =0,40 ) receberam a cor azul.
As questões objetivas do componente de Formação Geral, segundo o índice de facilidade, foram assim avaliadas : das oito questões, nenhuma teve o índice de facilidade classificado como muito fácil ou muito difícil.
Uma questão foi classificada como fácil, por ter índice de acertos situado na faixa entre 0,61 e 0,85 ( de 61,0 \% a 85,0 \% de acertos ).
Cinco questões foram consideradas de dificuldade média, situando -se no intervalo entre 0,41 e 0,60 do índice de facilidade, ou seja, houve entre 41,0 \% e 60,0 \% de acertos.
Duas questões foram classificadas como difícil, situando -se no intervalo entre 0,16 e 0,40.
Como já comentado, para análise das questões objetivas relativas à Formação Geral segundo o poder de discriminação, utilizou -se, o índice de discriminação ponto bisserial.
O Gráfico 3.13, para exemplificar, analisa o comportamento da questão de número 7 de Formação Geral.
Trata -se da segunda questão mais fácil e a que obteve o segundo maior índice de discriminação dessa parte da prova.
Neste gráfico, cada uma das cinco curvas representa o percentual de respostas em determinada alternativa da questão, em função da nota dos estudantes nesta parte da prova ( Formação Geral/Múltipla Escolha ), antes de possíveis eliminações pelo critério do ponto bisserial.
A curva em verde corresponde à alternativa B, a correta para esta questão.
Assim, observa -se que entre os estudantes com notas mais baixas, nessa parte do exame, a Nesta análise as questões foram assim avaliadas : seis das oito questões apresentaram índices acima de 0,40 e, assim, foram classificadas com índice muito bom para esse grupo de alunos ; uma questão teve índice de discriminação bom, entre 0,30 e 0,39.
Foi classificada com índice médio apenas uma das questões.
Nenhuma questão teve índice de discriminação fraco.
O índice de facilidade variou de 0,30 a 0,62, e o de discriminação, de 0,27 a 0,50.
As questões com índices de discriminação muito bom, de números 1, 2, 3, 4, 6 e 7, figuraram entre as mais fáceis desse conjunto : uma classificada na categoria fácil ( questão 1 ) do índice de facilidade e cinco na categoria médio ( questões 2, 3, 4, 6 e 7 ).
Em particular, a questão 2 foi a que apresentou maior poder discriminatório, com índice 0,50, e foi também a quinta mais fácil, com uma proporção de 0,53 de acertos.
A questão de número 5 apresentou índice de facilidade 0,30, ou seja, um quantitativo de 30,0 \% dos estudantes conseguiu resolvê -la, dentro do universo de participantes.
Além disso, seu índice de discriminação foi médio.
A questão 8 obteve índice de discriminação bom, 0,35, e seu índice de facilidade foi difícil.
O Gráfico 3.13, para exemplificar, analisa o comportamento da questão de número 7 de Formação Geral.
Trata -se da segunda questão mais fácil e a que obteve o segundo maior índice de discriminação dessa parte da prova.
Neste gráfico, cada uma das cinco curvas representa o percentual de respostas em determinada alternativa da questão, em função da nota dos estudantes nesta parte da prova ( Formação Geral/Múltipla Escolha ), antes de possíveis eliminações pelo critério do ponto bisserial.
A curva em verde corresponde à alternativa B, a correta para esta questão.
Assim, observa -se que entre os estudantes com notas mais baixas, nessa parte do exame, a A situação mais frequente foi a escolha de uma das alternativas incorretas : a alternativa C ( em preto ) ou A ( em azul ).
Na medida em que a nota aumenta, indicando desempenho melhor nesta parte da prova, aumenta concomitantemente a proporção de estudantes que selecionaram a alternativa correta B, atingindo 100 \% para os alunos com 8 acertos.
Essa análise permite verificar como a questão discriminou os grupos de desempenho, justificando o alto índice obtido na questão.
Os gráficos relativos às demais questões de Formação Geral constam do Anexo I. 3.2.2 Componente de Conhecimento Específico A Tabela 3.6 apresenta as estatísticas básicas em relação às questões objetivas do componente de Conhecimento Específico da prova, por Grande Região.
A média do Brasil deste componente foi de 41,9.
A menor média foi observada nas regiões Norte e Centro-Oeste ( 41,2 ) e a maior na região Sul ( 43,0 ).
O desvio padrão de todo o Brasil foi 15,3, sendo o menor desvio padrão encontrado na região Norte ( 14,2 ) e o maior na região Nordeste ( 15,8 ).
A mediana de todo o Brasil foi 40,0.
Em todas as regiões do Brasil, exceto para a região Sul ( 45,0 ), a mediana também foi 40,0.
A nota máxima da prova ( 100,0 ) foi obtida, nas questões objetivas do componente de Formação Geral, por pelo menos um aluno em cada região, exceto na região Norte ( 85,0 ) e na região Sul ( 90,0 ).
A nota mínima foi zero em todas as regiões, sem exceção.
A Tabela 3.7 apresenta os índices de facilidade e discriminação ( ponto bisserial ) das questões objetivas do componente de Conhecimento Específico.
Para facilitar a diferenciação das questões usaram -se as mesmas cores da Tabela 3.5 para as diferentes classificações dos índices de facilidade e de discriminação.
Dentre as questões objetivas da parte da prova relativa ao componente de Conhecimento Específico, a de número 26 foi anulada pela Comissão.
Desse modo, a classificação quanto ao índice de facilidade foi estabelecida com base em 26 das 27 questões.
A partir dos índices obtidos, pode -se concluir que a maioria das questões objetivas da prova foi considerada pelo menos difícil : das 26 questões, treze foram classificadas como difíceis e três como muito difíceis.
Não houve questão classificada como muito fácil, ao passo que duas foram tidas como fáceis, na faixa de 0,61 a 0,85 do índice de facilidade, e outras oito consideradas médias, entre 0,41 e 0,60.
Já quanto aos índices de discriminação das questões objetivas do componente de Conhecimento Específico da prova, tem -se como resultado a seguinte classificação : oito das 26 questões válidas foram consideradas como boas, enquanto duas delas tiveram índice de discriminação muito bom.
Assim, para dez das 26 questões, os índices de discriminação foram bons ou muito bons.
Dentre as demais, dez delas foram classificadas como médias e outras seis como fracas, sendo dezesseis, por conseguinte, a quantidade de questões nos dois patamares mais baixos de discriminação.
Constata -se, assim, que a prova – no que se refere -se ao componente de Conhecimento Específico – possuía capacidade fraca de discriminar entre aqueles que dominam ou não o conteúdo.
Dentre as questões que alcançaram os maiores índices de discriminação, as de números 24 e 25 classificadas com índice muito bom, situando -se no intervalo de índices maiores ou iguais a 0,40, também foram classificadas na categoria fácil.
A questão de número 31 foi a mais difícil dentre as 26 questões específicas válidas, com baixo índice de facilidade, apenas 8,0 \% de acertos.
Essa questão apresentou poder discriminatório igualmente baixo, 0,07, o que comprova ter sido esta a mais difícil para os estudantes.
Destaca -se, também, a questão 17, com índice de facilidade 0,10, o que, em termos percentuais, corresponde a 10,0 \% de estudantes que responderam acertadamente, obtendo, ainda, -0,02 de índice de discriminação.
Tais questões foram, portanto, duas das mais difíceis da prova.
As questões 31 e 17 foram eliminadas do cômputo da nota final pelo critério do ponto bisserial.
Além destas duas, as demais questões com índice fraco de discriminação, questões 10, 13, 23 e 32 também não foram computadas.
A título de exemplo das análises do comportamento das questões objetivas, o Gráfico 3.14 analisa a questão 24 do componente de Conhecimento Específico.
Esta foi a questão mais fácil da prova, apresentando índice de facilidade 0,81, ou seja, 81,0 \% dos estudantes assinalaram acertadamente a opção D, correspondente ao gabarito.
Seu índice de discriminação foi igual a 0,42, classificado como muito bom.
Neste gráfico, cada uma das cinco curvas representa o percentual de respostas em determinada alternativa da questão 24, em função da nota dos estudantes nesta parte da prova, antes de possíveis eliminações de questões pelo critério do ponto bisserial.
A alternativa correta D, representada no gráfico pela curva em roxo, foi escolhida em maiores proporções pelos alunos com desempenho melhor nesta parte da prova.
Já as alternativas incorretas, também denominadas distratores, foram selecionadas principalmente por aqueles com notas mais baixas.
Aqueles com nota zero, na sua quase totalidade deixaram esta questão em branco ou marcaram mais de uma alternativa, comportamento considerado inválido.
A proporção de alunos que selecionou a resposta correta D aumenta gradativamente, chegando a atingir 100 \% a partir de 20 acertos, enquanto a proporção dos que escolheram alternativas incorretas decai, a partir da primeira nota não nula, como função da nota nesta parte da prova.
Os gráficos relativos às demais questões do componente de Conhecimento Específico constam do Anexo I. 3.3 ANÁLISE DAS QUESTÕES DISCURSIVAS 3.3.1 Componente de Formação Geral As análises dos resultados de desempenho dos estudantes de Direito nas duas questões discursivas relativas à Formação Geral encontram -se na Tabela 3.8 e no Gráfico 3.15.
Na tabela 3.8 observa -se que as notas médias foram mais baixas nesse conjunto de questões do que no das objetivas.
Os estudantes de todo o Brasi0l obtiveram, em Formação Geral, média 49,3 nas questões objetivas e 41,1 nas questões discursivas.
A mediana também confirma o pior desempenho dos alunos de todo o Brasil nas questões discursivas do componente de Formação Geral.
Enquanto essa estatística foi de 45,0 para questões discursivas, para as questões objetivas essa estatística foi de 50,0.
Pode -se, também, notar um aumento do desvio padrão de 21,1, nas questões objetivas do componente de Formação Geral dos alunos de todo o Brasil, para 22,6 nas questões discursivas do mesmo componente.
Como já comentado, a mediana de todo o Brasil, neste componente, foi 45,0, assim como em todas as regiões, exceto na região Norte ( 47,5 ).
A nota máxima ( 97,5 ) foi obtida, nas questões discursivas do componente de Formação Geral, por pelo menos um aluno na região Sudeste.
As demais notas máximas foram 95,0 na região Nordeste e 92,5 nas regiões Norte, Sul e Centro-Oeste.
A nota mínima zero foi obtida em todas as regiões, sem exceção.
O Gráfico 3.15 representa a distribuição das notas nas questões discursivas no componente de Formação Geral.
A moda desta distribuição ocorre nos intervalos, ( 40 ; 50 ] e ( 50 ; 60 ].
Nota -se, também, uma grande concentração no intervalo [ 0 ; 10 ] devido à grande quantidade de notas zero e à alta frequência de alunos que deixaram este tipo de questão em branco.
A distribuição possui assimetria à esquerda, com coeficiente de assimetria -0,53.
Pelo gráfico pode -se observar uma maior concentração de notas acima do intervalo ( 40 ; 50 ].
Em todas as regiões o coeficiente de assimetria é negativo, o maior em módulo ocorreu na região Sudeste ( -0,58 ), e o menor, na região Nordeste ( -0,42 ).
Na sequência, os resultados verificados para cada uma das questões discursivas de Formação Geral serão apresentados, estabelecendo -se relações com os conteúdos abordados em cada uma delas.
Os comentários da Banca de docentes corretores a respeito do observado na correção das respostas dos estudantes, suas impressões e conclusões serão apresentados junto à análise de cada questão.
Cumpre esclarecer que, tendo em vista que as questões discursivas de Formação Geral são padronizadas, ou seja, constam de todas as provas, os comentários da Banca são 3.3.1.1 Análise da Questão Discursiva 1 do Componente de Formação Geral Os dados de Direito, obtidos a partir das respostas à questão 1, encontram -se na Tabela 3.9 e no Gráfico 3.16.
Nessa questão – de pior desempenho dentre as duas de Formação Geral – os alunos de todo Brasil tiveram média, 37,2.
A maior média para a questão 1 foi obtida na região Sudeste ( 38,2 ), e a menor, na região Centro-Oeste ( 35,6 ).
Quanto à variabilidade das notas, o desvio padrão de todo o Brasil foi 25,9.
O menor desvio padrão foi obtido na região Sudeste ( 25,3 ) e o maior desvio padrão foi obtido na região Nordeste ( 26,9 ).
As medianas e as notas mínimas da questão discursiva 1 foram as mesmas para todas as regiões do Brasil, respectivamente, 40,0 e 0,0.
As notas máximas foram 100,0 nas regiões Nordeste e Sudeste, e 95,0 nas regiões Norte, Sul e Centro-Oeste.
O Gráfico 3.16 mostra a distribuição das notas na questão discursiva 1 do componente de Formação Geral.
Observa -se que a maior frequência corresponde aos alunos que deixaram a questão em branco.
Considerando apenas os estudantes que tentaram resolver a questão, a moda ocorre no intervalo ( 40 ; 50 ].
Também observa -se grande concentração de notas no intervalo ( 30 ; 40 ].
O coeficiente de assimetria foi negativo ( -0,03 ), indicando uma assimetria à esquerda.
Nas regiões Norte ( -0,04 ), Sudeste ( -0,07 ) e Sul ( -0,02 ), o coeficiente de assimetria também foi negativo.
Já nas regiões Nordeste ( 0,02 ) e Centro-Oeste ( 0,04 ), o coeficiente de assimetria foi positivo, indicando uma assimetria à direita. 3.3.1.2 Comentários sobre a correção das respostas à Questão Discursiva 1 A questão aborda um tema bastante polêmico o qual vem sendo amplamente discutido, a conciliação entre o desenvolvimento econômico e a preservação dos recursos naturais.
Busca -se observar a relação entre o estímulo governamental à compra de veículos novos e os impactos ambientais decorrentes desta política pública.
Sua compreensão implica no conhecimento de uma temática atual e na reflexão critica a respeito das decisões governamentais em sua correlação com o cenário econômico internacional e os problemas e soluções de ordem ambiental.
Embora o comando da questão solicitasse explicitamente que a resposta fosse fornecida em um texto único ( “ redija um texto dissertativo sobre... ” ), foi significativamente alta a proporção de respondentes que atendeu a cada item ( a, b e c ) do enunciado de maneira isolada, como partes estanques da mesma resposta.
No que se segue encontra -se uma avaliação da correção de cada um dos três itens. a - “ conceito de desenvolvimento sustentável ” A variedade de respostas a esse item confirmou a importância do tema para avaliar a formação geral dos estudantes e cursos, já que ( i ) é atual ; ( ii ) foi abordado nos meios de comunicação ( especialmente no ano de 2012, em que o Brasil sediou a “ Rio+20 ” ) ; ( iii ) é relevante não apenas no Brasil como em todos os principais fóruns internacionais ; e ( iv ) é comum a diversas áreas de conhecimento.
Outros dois aspectos positivos merecem destaque : ( i ) a grande incidência de respostas que apontaram a relação entre gerações como aspecto essencial da ideia de sustentabilidade socioambiental – no mais, as respostas apresentaram pelo menos a tentativa de indicar uma definição de desenvolvimento sustentável ou a abordagem positiva dessa mesma noção, evidenciado compreensão do enunciado ; e ( ii ) o exercício crítico que a questão propunha, no sentido de exigir dos respondentes a compreensão das contradições de políticas públicas e a apresentação de propostas de soluções para conciliar o incentivo do crescimento econômico e a promoção da sustentabilidade socioambiental, a partir de uma perspectiva transformadora da realidade.
Apesar dos aspectos positivos apontados no parágrafo anterior, a maioria dos estudantes não alcançou o máximo da pontuação.
De forma geral, trataram o desenvolvimento sustentável como sinônimo de não agressão ambiental.
Alguns confundiram desenvolvimento sustentável com desenvolvimento econômico, e outros faziam referência apenas à preocupação com a conservação dos recursos mencionando a preocupação com gerações futuras.
Deve -se, ainda, destacar o significativo número de alunos que apresentou uma definição aparentemente copiada do enunciado de outra questão da prova, como por exemplo, a questão objetiva 31 da prova de Administração. b - “ conflito entre o estímulo à compra de veículos automotores e a promoção da sustentabilidade ” Esse foi o item identificado como aquele no qual as respostas dos estudantes avaliados mais se aproximaram da pontuação máxima prevista no padrão de resposta.
Ou seja, a ampla maioria dos respondentes foi bem-sucedida em apontar a relação entre o aumento da quantidade de veículos automotores em circulação e os danos socioambientais que isso acarreta no cenário urbano, devido ao aumento da emissão de gases poluentes.
No item b, as respostas corretas se concentraram em críticas ao incentivo oferecido pelo governo, referências ao problema dos congestionamentos e ao aumento da poluição.
Contudo, a maioria das respostas não explicitava a correlação entre o aumento da circulação de veículos automotores e o agravamento dos problemas ambientais.
Um número expressivo de alunos apresentou uma resposta em que era dado destaque à política de redução do IPI como um mecanismo para renovação da frota de automóveis ou para a aquisição do transporte próprio e da consequente redução da dependência ao transporte coletivo.
Também foram notados alguns tipos de padrões recorrentes nas tentativas de resposta a este item.
Dois desses padrões se destacaram : ( i ) as declarações genéricas de condenação do “ consumismo ” de uma “ nova classe média ” ou da busca pelo “ status ” que a posse de um automóvel particular proporcionaria, como se fosse possível atribuir a essa ordem de coisas a causa principal do aumento dos congestionamentos nas grandes cidades ; e ( ii ) a repetição da avaliação, presente no enunciado da questão, que aponta de forma superficial a contradição entre o estímulo à demanda por automóveis particulares e a promoção da sustentabilidade. c - “ ações de fomento ao transporte urbano sustentável no Brasil ” Este foi o item que ofereceu maior facilidade de aplicação do padrão de resposta, mas foram poucos os casos de pontuação integral.
Dentre as soluções mais destacadas pelos alunos apareceram : ( i ) o incentivo ao uso de energias alternativas ; e ( ii ) o estímulo ao transporte coletivo ( metrô, trem e ônibus ), além do uso de bicicletas.
Importante destacar o fato de que muitos apresentavam respostas com referências à necessidade de incentivo ao transporte urbano, nas quais este termo era utilizado como sinônimo de transporte coletivo.
A avaliação da questão 1 trouxe à tona problemas estruturais da educação básica, que se refletem no ensino superior.
Foi muito grande o número de provas em que o aluno demonstrou não compreender o enunciado da questão, apresentando respostas que fugiam ao tema proposto, mesmo quando bem redigidas.
Uma resposta recorrente fazia referência às consequências da ampla oferta de crédito à população e aos problemas relacionados à inadimplência, situação que, embora, verdadeira, não estava relacionada à questão proposta.
Houve, ainda, os casos em que as respostas apresentadas não só se afastaram do tema como não reuniam qualquer tipo de argumentação coerente.
E, por fim, destaca -se a existência de respostas incoerentes e desconexas.
Outro aspecto geral diz respeito ao grande número de provas em branco ou em que o aluno apresentou a resposta a outro item no espaço da questão 1.
Considerando -se o objetivo geral deste exame, avaliar a qualidade dos cursos de graduação oferecidos no país e contribuir para a melhoria da qualidade do ensino superior, destaca -se a imensa responsabilidade do processo que envolve a elaboração deste diagnóstico.
Nesse sentido, é urgente alertar para a ocorrência de um número significativo de respostas que se afastam de um padrão mediano de qualidade esperado de um aluno concluinte em um curso de graduação, especialmente em se considerando o fato de que a questão proposta faz referência a tema intensamente explorado pela mídia.
Outro aspecto geral diz respeito ao grande número de provas em branco ou em que o aluno apresentou a resposta a outro item no espaço da questão 1.
Considerando -se o objetivo geral deste exame, avaliar a qualidade dos cursos de graduação oferecidos no país e contribuir para a melhoria da qualidade do ensino superior, destaca -se a imensa responsabilidade do processo que envolve a elaboração deste diagnóstico.
Nesse sentido, é urgente alertar para a ocorrência de um número significativo de respostas que se afastam de um padrão mediano de qualidade esperado de um aluno concluinte em um curso de graduação, especialmente em se considerando o fato de que a questão proposta faz referência a tema intensamente explorado pela mídia. 3.3.1.3 Análise da Questão Discursiva 2 do Componente de Formação Geral A Tabela 3.10 mostra que o desempenho dos estudantes na questão 2 ( média 45,0 ) foi superior ao obtido na questão de número 1 ( média 37,2 ).
A região Norte foi aquela onde a média, nessa questão, foi maior ( 47,4 ), e a de menor média foi a região Nordeste ( 43,0 ).
Quanto à variabilidade das notas, o desvio padrão de todo o Brasil foi 25,5, inferior ao obtido na questão de número 1 ( 25,9 ).
O maior desvio nessa questão foi obtido na região Nordeste ( 27,2 ), enquanto o menor foi obtido na região Norte ( 24,6 ).
As medianas e as notas mínimas da questão discursiva 2 foram as mesmas para todas as regiões do Brasil, respectivamente, 50,0 e 0,0.
A nota máxima ( 100,0 ) foi obtida por pelo menos um aluno na região Centro-Oeste.
A nota máxima ( 95,0 ) foi obtida nas demais regiões ( Norte, Nordeste, Sudeste e Sul ).
O Gráfico 3.17 mostra a distribuição das notas na questão discursiva 2 do componente de Formação Geral.
Neste gráfico, observa -se a moda no intervalo ( 50 ; 60 ].
Destacam -se, também, a concentração no intervalo ( 40 ; 50 ], próximo à modal, e as notas em branco.
O coeficiente de assimetria das notas nesta questão é negativo ( -0,64 ), indicando que a distribuição é assimétrica à esquerda.
O mesmo ocorre em todas as regiões, variando entre -0,74 na região Norte e -0,50 na região Nordeste. 3.3.1.4 Comentários sobre a correção das respostas à Questão Discursiva 2 O conteúdo da questão era bastante adequado às Diretrizes do INEP para o ENADE, além de atual e de abrangente.
Porém, exigia -se a articulação de muitas variáveis.
Buscando aferir, a partir da apresentação de uma charge e de um texto introdutório, as competências e habilidades relacionadas à leitura, à interpretação, à análise crítica, à capacidade argumentativa, à elaboração de sínteses e de avaliação prospectiva sobre tema da atualidade, a questão tratava da violência e suas repercussões no contexto atual.
O comando da questão solicitava, a partir da definição de violência da OMS e de duas imagens, que o aluno elaborasse um texto, de caráter dissertativo, sobre a “ violência na atualidade ”.
Embora tal solicitação fosse explícita ( “ redija um texto dissertativo sobre... ” ), como na questão 1, foi alta a proporção de alunos que respondeu a cada item do enunciado isoladamente, como partes independentes.
Verificaram -se poucas respostas em branco, como também foram poucas as respostas contendo protestos contra o ENADE, muito frequentes em anos anteriores, o que significa um avanço para a afirmação do ENADE como um exame aceito socialmente. a -"tecnologia e violência"Tanto as charges quanto o texto apresentado no enunciado induziram os alunos a focarem suas respostas nos meios de comunicação.
Além disso, algumas respostas estabeleciam a tecnologia, em si, como causa da violência tanto nas escolas ( em particular ), quanto na sociedade ( no geral ).
Muitos respondentes tiveram dificuldade para fazer a relação entre tecnologia, violência e a escola, de modo adequado, centrando -se na violência nas escolas.
Além disso, a grande maioria das respostas girava em torno da exposição da violência através da TV, internet e jogos digitais, sendo que predominaram as abordagens com ênfase em programas e filmes veiculados pela televisão.
Poucos foram os estudantes que estabeleceram relação entre o contexto socioeconômico e a violência. b -"causas e consequências da violência na escola"Embora a solicitação desse item fosse objetiva e atendida por grande parte dos alunos, inúmeras vezes, observou -se confusão entre o significado de causa e de consequência.
Em outras inúmeras vezes, as respostas apontavam apenas para uma das classes, causas ou consequências.
Foram frequentes as respostas que atribuíam o aumento da violência à ineficiência do poder público, apontando a falta de investimentos, o baixo salário dos professores, a ausência de profissionais qualificados como causas.
Foram frequentes também os que apontavam a violência doméstica, o bullying, e a transferência do imaginário dos filmes e jogos para a realidade, como causas da violência. c -"proposta de solução para o problema da violência na escola"Sem dúvida, há uma grande diversidade de abordagens possíveis para apresentação de uma proposta de solução para o problema da violência, o que o padrão de resposta buscou contemplar.
No entanto, houve grande número de respostas baseadas em opiniões pessoais dos estudantes e vinculadas ao senso comum.
As respostas giraram em torno da necessidade de mais investimentos na educação por parte do poder público e seu rebatimento no maior aparelhamento das escolas, salário dos professores etc.
Além disso, foram muitas as respostas que trataram da necessidade de maior conscientização dos pais, alunos e professores com a realização de um trabalho conjunto que envolvesse todos os segmentos da sociedade.
De um modo geral, os estudantes demonstraram dificuldades em relação ao nível de abrangência exigido, falta de conteúdo, inabilidade na elaboração de texto escrito.
Da grande quantidade de provas corrigidas, o que ficou evidente, foi a dificuldade de estabelecer as conexões requeridas na questão 2.
A violência nas escolas foi bastante associada ao reflexo do meio onde a criança e o adolescente mantém os seus vínculos afetivos.
Surgiram muitas respostas revelando críticas à intolerância no núcleo familiar, à jornada de trabalho dos pais, ao distanciamento da mãe nas atividades domésticas etc.
Em boa parte das respostas, os estudantes lançaram a causa da violência escolar exclusivamente na ausência de participação familiar na vida dos filhos, às vezes criando uma cisão com a proposta da questão, já que abandonavam totalmente a temática da tecnologia em suas reflexões.
Além da ausência dos pais no cotidiano da escola, divórcios e atos de violência contra crianças foram considerados fatores que delegariam para a escola novas funções.
Também foram mencionadas a má formação de professores e a insatisfação com o ambiente escolar.
Em inúmeros casos, ficou evidente o anseio de se fazer da escola um lugar de resgate tanto da criança e do adolescente quanto da família, compreendendo -se a escola em sentido social e assistencial.
Segundo este amplo grupo de concluintes, a escola deveria oferecer -se como centro de cultura e lazer, como espaço para o acompanhamento psicológico e serviço social às comunidades diretamente envolvidas.
Ressalte -se o constante aparecimento das sugestões de horário integral, melhores salários para professores e iniciativas para melhoria na formação de professores.
As propostas de solução para acabar com a violência na escola se detiveram em aspectos muito gerais e poucos alunos detalharam uma proposta de intervenção, apresentando ações concretas.
Dentre as propostas de solução surgiram : educação familiar, ensino religioso, maior segurança nas escolas, com utilização de aparelhos para detecção de armas, e o retorno do papel do professor como um profissional respeitado pelos poderes municipais e estaduais.
Um percentual significativo de provas apresentou resposta simplista, com ênfase em determinados pontos como : religião ( sem demonstração do conceito de laicidade ), políticas sociais ( sem demonstração do domínio de conceitos que envolvem Estado, Governo e Foi grande a incidência de respostas que se referiam às leis, à vigilância e maior controle dos pais e professores ; no entanto, não demonstram entender a noção de legalidade para os diferentes segmentos sociais e como os mesmos se encontram submetidos a diferentes formas de poder.
Alguns alunos, ainda que raros, trataram a questão com a profundidade esperada de concluintes de cursos universitários.
Por outro lado, foram vários os textos que apresentaram uma argumentação simplória e desconexa, com forte visão paternalista sobre o papel do governo e do professor, como se tais instâncias pudessem intervir na educação familiar.
Enquanto alguns conseguiram desenvolver o tema proposto com elegância e desenvoltura, encadeando o tema de forma harmônica, com excelente nível de linguagem escrita e correção gramatical, outros sequer conseguiram conectar as ideias e construir um texto razoável, utilizando, inclusive, um linguajar mais propício à linguagem falada do que à escrita.
Cabe também ressaltar a quantidade de respostas carregadas de conteúdo com juízo de valor discriminatório ou extremamente conservador, do tipo : “ mulher não deve trabalhar fora ” ; “ amar e temer a Deus é a solução ” ; obrigatoriedade do ensino religioso e de uma doutrina religiosa ; imposição de uma educação do tipo militar nas escolas ; censura nos meios de comunicação etc.
Outras soluções apontadas para resolver o problema da violência nas escolas estavam associadas à, por exemplo : redução da maioridade penal, defesa da pena de morte, milícias armadas na porta das escolas, punições como o retorno da palmatória.
O Estatuto da Criança e do Adolescente foi uma peça legal criticada, e recebeu, em diversas respostas, o título de responsável pelo sentimento de impunidade de adolescentes infratores.
Contudo, algumas dessas proposições, desde que bem contextualizadas e formuladas com boa argumentação, foram pontuadas de forma a respeitar o direito democrático de pluralidade opinativa.
É imperioso ressaltar que a falta de inteligibilidade segue representando um aspecto merecedor de providências na formação universitária, uma vez que exprime uma habilidade importante para qualquer profissional de nível superior.
Além de terem sido poucas respostas boas e claras, foi possível identificar a deficiência de argumentação, além dos inúmeros problemas gramaticais ( ortografia, concordância, pontuação, regência, etc. ).
Salienta -se, nesse sentido, o uso indiscriminado de letras maiúsculas e minúsculas numa mesma frase, sem obedecer ao critério formal da língua portuguesa. 3.3.2 Componente de Conhecimento Específico Na parte da prova relativa às questões discursivas no componente de Conhecimento Específico ( Tabela 3.11 ), observa -se que a média foi mais baixa do que para as questões discursivas do componente de Formação Geral.
Enquanto no componente de Formação Geral a média das questões discursivas para estudantes de Direito de todo o Brasil foi 41,1, na parte de Conhecimento Específico a média foi 22,7.
A maior média deste componente foi obtida pelos estudantes da região Sudeste ( 24,1 ), e a menor, pelos da região Centro-Oeste ( 20,9 ).
Quanto à variabilidade das notas, o desvio padrão de todo o Brasil foi 18,0.
O maior desvio padrão foi encontrado na região Nordeste ( 18,4 ), e o menor, na região Sul ( 17,7 ).
A maior nota máxima, nas questões objetivas do componente de Conhecimento Específico, foi obtida na região Sul ( 98,3 ) enquanto a menor nota máxima foi encontrada nas regiões Sudeste e Centro-Oeste ( 86,7 ).
Além disso, a mesma nota mínima ( 0,0 ) foi obtida por alunos de todas as regiões do Brasil sem exceção.
A mediana do Brasil como um todo foi 21,7, também encontrada nas regiões Norte e Sul.
As demais medianas foram : 23,3 na região Sudeste, 20,0 na região Nordeste e 18,3 na região Centro-Oeste.
O Gráfico 3.18 representa a distribuição das notas nas questões discursivas no componente de Conhecimento Específico.
Essa distribuição tem moda no intervalo de [ 0 ; 10 ].
Nota -se que a partir desse intervalo há outra concentração de notas no intervalo ( 20 ; 30 ].
O coeficiente de assimetria é positivo ( 0,38 ), indicando uma assimetria à direita.
Este coeficiente também é positivo em todas as regiões do Brasil ( entre 0,28 na região Sudeste e 0,50 na região Centro-Oeste ). 3.3.2.1 Análise da Questão Discursiva 3 do Componente de Conhecimento Específico Na questão 3, cujos resultados aferidos encontram -se descritos na Tabela 3.12, a média dos estudantes de todo o Brasil foi 24,1.
A menor média nessa questão foi obtida pelos alunos da região Centro-Oeste ( 22,1 ), enquanto a maior média foi obtida na região Sudeste ( 25,7 ).
Quanto à variabilidade das notas, o desvio padrão de todo o Brasil foi 23,3.
O maior desvio padrão foi obtido na região Nordeste ( 23,6 ), enquanto o menor foi obtido na região Centro-Oeste ( 22,7 ).
A mesma nota máxima ( 100,0 ), mediana ( 20,0 ) e nota mínima ( 0,0 ) foram alcançadas em todas as regiões do Brasil, sem exceção.
O Gráfico 3.19 mostra a distribuição das notas na questão discursiva 3, do componente de Conhecimento Específico.
Essa distribuição é bimodal, com modas nas questões em branco e no intervalo ( 20 ; 30 ].
O coeficiente de assimetria é positivo para o Brasil como um todo ( 0,74 ), e também para todas as regiões ( entre 0,67 para a região Sudeste e 0,85 para a região Nordeste ). 3.3.2.2 Comentários sobre as respostas à Questão Discursiva 3 O enunciado era claro, elaborado com linguagem técnica, abordando conteúdo inserido nas Diretrizes Gerais, com nível médio de profundidade.
A divisão das perguntas também ficou clara indicando, passo a passo, o caminho da resposta a ser seguido que deveria ter por base a legislação indicada no texto.
Verificou -se uma dificuldade dos examinandos quanto à análise da licitude do objeto e de apresentar as formalidades do contrato tendo em vista que a maior parte considerou o terceiro item prejudicado pela resposta formulada aos dois primeiros.
As respostas fornecidas pelos estudantes abordaram a questão da capacidade civil e para o trabalho como previstas na nossa legislação constitucional, civil e trabalhista, indicando, no mais das vezes, a própria lei referida no enunciado.
Muitos não conseguiram diferenciar a questão da capacidade com a do objeto licito do contrato.
Nesse caso, inúmeros examinandos aduziram a nulidade da avença diante da menoridade indicada no enunciado.
No pertinente ao último item também houve dificuldade, na maioria das vezes, de identificar os elementos do contrato, pois a resposta estaria prejudicada ao deduzir a incapacidade de contratar e a nulidade do contrato.
Muitas respostas apresentaram o tema de menor aprendiz, o que foi aproveitado.
Em relação ao item a, ao revelar a incapacidade do agente menor em celebrar a avença, e indicando a norma aplicável, seja da Constituição, do Código Civil ou da Lei indicada no texto do enunciado da questão, o estudante obteve alguma pontuação.
Os que responderam afirmativamente à incapacidade em contratar, indicando a norma constitucional aplicável, a norma do Código Civil, a norma da Consolidação das Leis do trabalho ou, ainda, a norma citada no texto, receberam notas mais elevadas.
A seguir o estudante concluiu ser o objeto lícito por inexistir vedação legal quanto ao contrato.
As respostas que obtiveram as maiores notas, tendo em vista que atenderam aos objetivos da questão formulada, variaram da descrição detalhada dos requisitos legais, indicado na lei prevista no texto, à mera remissão ao artigo que continha os requisitos.
As respostas consideradas medianas foram incompletas ou por não indicarem a legislação aplicável ( Código Civil, CLT, Constituição e/ou lei especial ) ou por não responderem de maneira adequada aos três itens perguntados.
Já as respostas com as notas mais baixas foram as que afirmaram a existência de capacidade, mesmo diante do texto da lei, ou apresentaram respostas completamente dissonantes, inclusive indicando como objeto ilícito sempre que há incapacidade.
Na mesma linha, aquelas que apontaram nulidade do contrato firmado por menor incapaz, ao invés de indicar que o objeto seria lícito, mas haveria incapacidade.
Também foram equivocadas as respostas que afirmaram que seria a possibilidade estabelecer contrato com representação dos pais, colidindo com as normas legais.
Algumas respostas defenderam que os menores deveriam estudar ao invés de trabalhar, mesmo em atividades equiparadas ao esporte.
Aduziram outros que a atividade seria perigosa e, portanto, vedada à inclusão de menores, mesmo como aprendizes. 3.3.2.3 Análise da Questão Discursiva 4 do Componente de Conhecimento Específico A Tabela 3.13 contém as informações relativas à questão 4 do conjunto de questões do componente de Conhecimento Específico.
O desempenho dos estudantes de todo o Brasil nesta questão foi aproximadamente igual ao desempenho na questão de número 3.
A média geral do Brasil foi 24,1, sendo a menor média registrada na região Centro-Oeste ( 21,5 ) e a maior na região Sudeste ( 26,2 ).
A mesma nota máxima ( 100,0 ), assim como a mesma nota mínima ( 0,0 ), foi atingida em todas as regiões, sem exceção.
A mediana para o Brasil como um todo foi 5,0, a mesma obtida na região Sul.
As demais medianas foram 15,0 na região Sudeste e 0,0 nas regiões Norte, Nordeste e Centro-Oeste.
O Gráfico 3.20 representa a distribuição de notas na questão discursiva 4, no componente de Conhecimento Específico.
Como nas outras questões discursivas, o número de estudantes que deixou a questão em branco foi muito grande, correspondendo à moda da distribuição.
Dentre as notas dos que resolveram a questão 4, o intervalo modal foi ( 50 ; 60 ], e vale destacar o intervalo ( 60 ; 70 ] que foi muito próximo ao seu intervalo anterior.
Os coeficientes de assimetria são sempre positivos, tanto na distribuição como um todo ( 0,68 ) quanto na análise dos resultados dos alunos por região ( entre 0,54 para a região Sudeste e 0,85 para as regiões Nordeste e Centro-Oeste ). 3.3.2.4 Comentários sobre as respostas à Questão Discursiva 4 O enunciado era claro, elaborado com linguagem técnica, com temática inserida nas Diretrizes Gerais e nível médio de profundidade.
A divisão das perguntas também ficou clara indicando, passo a passo do caminho da resposta a ser seguido.
Os estudantes não conseguiram apreender, em sua ampla maioria, as consequências do fenômeno da atribuição vinculando -o à sanção decorrente.
As respostas foram genéricas e basearam -se no próprio texto do enunciado sem aditar nova informação.
Paradoxalmente, ao exemplificar as relações jurídicas, as respostas apresentadas foram ao encontro do padrão de resposta, indicando exemplos satisfatórios em relação aos objetivos perseguidos pelo padrão de resposta.
O exemplo mais citado foi a relação contratual. objetivos perseguidos pelo padrão de resposta.
O exemplo mais citado foi a relação contratual.
A correção buscou pontuar as respostas mais próximas do padrão apresentado.
Assim, ao indicar os elementos do direito subjetivo, do dever jurídico, da relação jurídica e exemplificando com o contrato de compra e venda, como a maioria dos estudantes apresentou, houve pontuação parcial dos três itens perguntados.
Essa questão, por estar relacionada à disciplina vinculada à Teoria Geral do Direito, foi a que contou com maior dificuldade de desenvolvimento por parte dos examinandos.
Muitos não apreenderam o significado dos conceitos enunciados no texto apresentado no enunciado.
Boas respostas foram aquelas que indicaram a fundamentação e o exemplo pretendido que foi, na maior parte das vezes, o contrato de compra e venda.
Houve referência também ao contrato de locação, ao contrato de prestação de serviços, ao casamento, ao divórcio, e às relações processuais.
Foram consideradas medianas as respostas incompletas que não indicaram os institutos aplicáveis, quer no aspecto do direito subjetivo/dever jurídico, quer no aspecto da atribuição ( sanção ), fugindo parcialmente do padrão de resposta.
Como exemplo de respostas fracas estão aquelas que : a ) não consideraram a bilateralidade das relações afirmando que existe uma opinião que fundamenta o direito subjetivo ; b ) indicaram uma responsabilidade sem relação ; c ) indicaram o casamento como relacionado diretamente a uma garantia jurídica. 3.3.2.5 Análise da Questão Discursiva 5 do Componente de Conhecimento Específico A Tabela 3.14 contém as informações relativas à questão 5 do conjunto do componente de Conhecimento Específico.
O desempenho dos estudantes nessa questão foi inferior ao das questões 3 e 4.
A nota média dos estudantes de todo o Brasil foi 19,7.
A maior média foi registrada na região Norte ( 20,9 ), enquanto a menor média foi registrada na região Sul ( 18,9 ).
Quanto à variabilidade das notas, o desvio padrão dos alunos do Brasil, como um todo, foi 20,3.
Enquanto o maior desvio foi encontrado na região Norte ( 21,4 ), o menor foi encontrado na região Sudeste ( 19,5 ).
A mesma nota máxima ( 100,0 ) e a mesma nota mínima ( 0,0 ) foram alcançadas em todas as regiões do Brasil, sem exceção.
A mediana para o Brasil com um todo foi 15,0, a A mesma nota máxima ( 100,0 ) e a mesma nota mínima ( 0,0 ) foram alcançadas em todas as regiões do Brasil, sem exceção.
A mediana para o Brasil com um todo foi 15,0, a mesma obtida nas regiões Nordeste, Sul e Centro-Oeste.
Já nas regiões Norte e Sudeste a mediana foi 20,0 na questão.
O Gráfico 3.21 representa a distribuição das notas da questão discursiva 5 do componente de Conhecimento Específico.
Destaca -se, como nas demais questões discursivas, o grande número de estudantes que deixaram esta questão em branco, correspondendo à moda.
Dentre as notas dos que resolveram a questão 5, o intervalo modal foi ( 20 ; 30 ].
Entre as três questões discursivas do Componente de Conhecimento Específico, esta foi a que apresentou o menor número de questões em branco.
O coeficiente de assimetria para todos os alunos é 1,31, indicando uma assimetria à direita, e para todas as regiões se mantém positivo, variando entre 1,28 para a região Norte e 1,34 para a região Centro-Oeste. 3.3.2.6 Comentários sobre as respostas à Questão Discursiva 5 A questão foi considerada de nível médio de profundidade e, como nas anteriores, a divisão das perguntas contribuiu para indicar o caminho da resposta a ser seguido.
A maioria dos estudantes não conseguiu apreender a excepcionalidade da situação.
Muitos identificaram corretamente a norma aplicável, de natureza constitucional, albergando o direito à inviolabilidade de domicilio, pretendida pelo padrão.
No entanto, não conseguiram entender que essa seria uma exceção pugnando pela ilegalidade da atuação dos policiais e alguns insistindo na necessidade de mandado judicial.
Também foram identificadas outras expressões como inviolabilidade da casa, da residência etc.
Muitos não indicaram o fundamento constitucional exigido pela resposta.
Diferentemente das questões anteriores, a pontuação das respostas como um todo foi dificultada pelo natural encadeamento exigido para a resposta.
Para a pontuação do item ( a ) ser integral, bastou a menção à proteção constitucional relativa à inviolabilidade de domicílio.
Também foram poucos os que identificaram a atuação dos policiais como estrito cumprimento do dever legal, mesmo aqueles que identificaram a atuação como legítima na trilha do padrão apresentado.
Verificou -se que os estudantes têm uma forte formação garantística, não acolhendo a atuação da polícia com facilidade, identificando o Estado como opressor e outorgando maior poder ao cidadão para se defender daquilo que muitos consideraram uma violação ao sistema jurídico que seria a invasão do domicilio sem mandado judicial, por ter sido a diligência praticada fora do horário normal, considerado diurno.
Como exemplos de respostas boas estão aquelas que preencheram o padrão apresentado, indicando a fundamentação e o exemplo pretendido que foi alcançado pela minoria dos estudantes que não conseguiram apreender que a hipótese do enunciado albergava uma exceção à proteção constitucional.
Por outro lado, também não houve a identificação da legitimidade da atuação dos policiais com o exercício do estrito cumprimento de dever legal.
Foram consideradas medianas as respostas que indicavam a norma constitucional, mas não chegavam à conclusão correta, ou as que não indicaram os institutos aplicáveis fugindo parcialmente do padrão de resposta apresentado.
Dentre as respostas consideradas fracas estão aquelas que : a ) não consideraram a existência de flagrante ; b ) consideraram ofendida a norma constitucional que protege a inviolabilidade de domicílio ; c ) indicaram a necessidade de mandado judicial para adentrar no domicílio dos indivíduos ; d ) consideraram a atividade dos policiais ilícita e com abuso de poder.

% matched lemmas: 
\end{textsample}
